%%-----------------------%%
%% Homomorphic Expansion %%
%%-----------------------%%
\section{RLWE-to-RGSW Expansion}

\subsection{ORAM}
\ifnotes
    \begin{itemize}
        \item We want to use RGSW to avoid relinearization
        \item We want to send packed ciphertexts to save on communication
        \item Therefore, we send RLWE and convert to RGSW on the server
        \item (Requires a single RGSW encryption $\mathsf{RGSW}(-s)$ to be sent as well)
        \item 1 RLWE for the index bits = size $\ell$
        \item $\ell$ RGSW's for the index bits = size $2\ell \times 2$ or smth
        \item Expanding = size $2\ell$ or something
    \end{itemize}
\fi
As established in \autoref{sec:rgsw}, \gls{rgsw} ciphertexts are enticing because the external product allows us to perform (binary) homomorphic multiplication without relinearization. However, the size of \gls{rgsw} makes the communication costs very high in a communication protocol. We wish to strike a balance between the cheap communication cost of packed \gls{rlwe} and the fast multiplication of \gls{rgsw}. The solution is to send \gls{rlwe} ciphertexts from the client, and \emph{expand} them homomorphically on the server.

Solution from \autocite{chillotti:onionring:2019}. \textbf{HomExpand:} Expands an \gls{rlwe} ciphertext into an \gls{rgsw} ciphertext.

This expansion is described in Algorithm 4 of \cite{chillotti:onionring:2019}, which builds on Algorithms 3 and 8 from the same paper. The algorithms are listed below (in slightly different notation) as \autoref{alg:ORAM:4}, \autoref{alg:ORAM:3}, and \autoref{alg:ORAM:8}, respectively.

%%%
%%% ORAM ALGORITHM 4 - HomExpand (RLWEtoRGSW)
%%%
\algnewcommand\algorithmicforeach{\textbf{for each}}
\algdef{S}[FOR]{ForEach}[1]{\algorithmicforeach\ #1\ \algorithmicdo}

%%%
%%% ORAM Algorithm 4 - HomExpand (RLWEtoRGSW)
%%%
\begin{algorithm}[ht!]
\caption{HomExpand}
\label{alg:ORAM:4}
\begin{algorithmic}[1]
    \State \textbf{Input:} $\mathbf{c}_k$ = RLWE$(\sum_{i=0}^{n-1}\frac{b_i}{nB^{k+1}}),\; b_i\in\{0,1\},\; 0 \leq k < \ell$
    %, \Statex $\quad\quad\quad\;\; k=0,1,\dots,\ell - 1$
    \State \textbf{Input:} $A = \mathrm{RGSW}(-s)$
    \State \textbf{Output:} $\mathbf{C}_i = \mathrm{RGSW}(b_i),\; 0 \leq i < n$
    \ForEach{$k$} % in $0, \dots, \ell - 1$}
        \State $\mathbf{c'}_k:= \mathsf{expandRlwe}(c_k)$
    \EndFor
    \For{$i = 0,\; i < n$}
        \State Initialize $\mathbf{C}_i$ as an empty $2\ell \times 2$ matrix of polynomials
        \For{$k = 0; k < \ell$}
            \State $\mathbf{C}_i[k] = A \boxdot \mathbf{c'}_k[i]$
            \State $\mathbf{C}_i[k + \ell] = \mathbf{c'}_k[i]$
        \EndFor
    \EndFor
    \State\Return $\{\mathbf{C}_j\}^{n-1}_{j=0}$
\end{algorithmic}
\end{algorithm}

\autoref{alg:ORAM:4} takes as input $\ell$ RLWE ciphertexts, which in \gls{spar} correspond to the $\ell$ encrypted bits of the index \textcolor{red}{$a_j$}. It also takes an RGSW ciphertext of the encrypted secret key % \textcolor{red}{(which key, automorphism? Sent by the client or generated by server?)}.

It begins by running a subroutine $\mathsf{expandRlwe}$ (\autoref{alg:ORAM:3}) which converts each of the $\ell$ RLWE ciphertexts into a matrix, like a \emph{partial} RGSW. Then these partial RGSW's are combined into a single (complete) RGSW.

%%%
%%% ORAM ALGORITHM 3 - RLWE Expansion subroutine ExpandRLWE
%%%
\input{Algorithms/oram-3}

This subroutine converts an RLWE ciphertext with $\ell$ slots into $\ell$ ciphertexts, each encoding a unique slot of the input. \textcolor{red}{Subs} is a function that... (EvalAutomorphism!)

Their method is fast for TFHE, but for our use-case, where $n << N$, it's faster to use the method described in \cite{cryptoeprint:2018/244} to repeatedly ``rotate'' the ciphertext $\ell$ times as in \autoref{alg:expand:bgv}. The FastRotate operation is supported by the OpenFHE as $\mathsf{EvalFastRotate}$.

Note that the components being scaled by $n$ is a side-effect of their implementation, and not required when using FastRotate, so the output is a list of $c_i =\mathsf{RLWE}(b_i)$.

%%
%% ORAM ALGORITHM 3 - MODIFIED
%%
%%%
%%% ORAM Algorithm 3 - RLWE Expansion subroutine ExpandRLWE
%%%
\begin{algorithm}[ht!]
\caption{expandRlwe}
\label{alg:expand:bgv}
\begin{algorithmic}[1]
    \State \textbf{Input:} $\mathbf{c} = \mathrm{RLWE}(\sum_{i=0}^{n-1} b_i X^i),\;b_i\in\{0,1\}$
    \State \textbf{Output:} $\mathbf{c}_i = \mathrm{RLWE}(b_i),\; 0 \leq i < n$
    \State $\mathbf{c}_0 := \mathbf{c}$
    \For{$i = 0,\; i \leq \log{n}$}
        \State do something...
        % \State $k = n/2^{i - 1} + 1$
        % \For{$b = 0; b < 2^i$}
        %     \State $\mathbf{c}_{2b} = c_b + \mathrm{Subs}(c_b, k)$
        %     \State $\mathbf{c}_{2b + 1} = c_b - \mathrm{Subs}(c_b, k)$
        % \EndFor
    \EndFor
    \State\Return $\{\mathbf{c}_j\}^{n-1}_{j=0}$
\end{algorithmic}
\end{algorithm}



\ifnotes
    \subsubsection{Substitution}
    Finally, the substitution is a bit more complex... Uses an automorphism key to rotate? Or something?
    
    %%% -----------------------
    %%% ORAM ALGORITHM 8 - Subs
    %%% -----------------------
    %%%
%%% ORAM Algorithm 8 - Subs
%%%
\begin{algorithm}
\caption{RLWE Substitution Algorithm $\mathsf{Subs}(c, k)$}
\label{alg:ORAM:8}
\begin{algorithmic}[1] % [1] adds line numbering every line
    \State \textbf{Input:} $\mathbf{c} = \mathrm{RLWE}_s(\mu(X))$; odd integer $k$
    \State \textbf{Input:} $KS_{s(X^k)\to s}[j]$ for $j = 1 : \ell_{ks}$
    \State \textbf{Output:} $\mathbf{c'} = \mathrm{RLWE}_s(\mu(X^k))$
    \State $\mathbf{c} = (c_0, c_1)$
    \State Let $\mathbf{c'} = (c'_0, c'_1) = (c_0(X^k), c_1(X^k))$
    \State Decompose $c'_0$ as $\sum_{j=1}^{\ell_{ks}} c'_0[j]/B^j_{ks} + \epsilon_{dec}$ with
    \Statex $\quad||c'_0[j]||_\infty \leq B_{ks}/2$ and $||\epsilon_{dec}||_\infty\leq 1/(2B^\ell_{ks})$
    \State\Return $\mathbf{c'} = (0, c_1) - \sum_{j=1}^{\ell_{ks}} c_0[j] \cdot KS_{s(X^k)\to s}[j]$
\end{algorithmic}
\end{algorithm}
\fi


\ifnotes
    %%----------------%%
    %%                %%
    %%----------------%%
    \subsection{Homomorphic Field Trace}
    New paper!
\fi